% PAGES: 272-272

\noindent\textbf{Minimum variance portfolio using tangency portfolio:} If
\[
w_T \;=\; \frac{1}{\bfone^t\Sigma_R^{-1}\bar\mu}\Sigma_R^{-1}\bar\mu
\]
is the weight of the $n$ assets in the tangency portfolio, then the asset weights
vector $w_{min}$ and the weight $w_{min,cash}$ of the cash position given by (9.13) and
(9.14) can also be computed as follows:
\begin{equation}
w_{min,cash} \;=\; 1 - \frac{\mu_P-r_f}{\bar\mu^tw_T};
\end{equation}
\begin{equation}
w_{min} \;=\; (1-w_{min,cash})w_T.
\end{equation}

The standard deviation of the return of the minimum variance portfolio is
\[
\sigma_{min} \;=\; \sqrt{w_{min}^t\Sigma_Rw_{min}}.
\]

The asset allocation for the minimum variance portfolio computed using the tangency
portfolio can be found using the pseudocode from Table 9.2.

\begin{table}[H]
\centering
\caption{Asset allocation of minimum variance portfolio from tangency portfolio}
\fbox{\begin{minipage}{0.85\linewidth}
Input:\\
$\Sigma_R = n\times n$ covariance matrix of the returns of $n$ assets\\
$\mu = n\times 1$ vector of expected values of the returns of $n$ assets\\
$r_f = $ risk--free rate\\
$\mu_P = $ required expected return of the portfolio
\bigskip

Output:\\
$w_{min} = $ asset weights vector for the minimum variance portfolio\\
$w_{min,cash} = $ weight of cash position for the minimum variance portfolio\\
$\sigma_{min} = $ standard deviation of the return of the minimum variance portfolio
\bigskip

$\bar\mu = \mu - r_f\bfone$\\
$x = \text{linear\_solve\_cholesky}(\Sigma_R,\bar\mu)$ \hfill\textit{// compute $x =
\Sigma_R^{-1}\bar\mu$}\\
$w_T = \dfrac{1}{\bfone^tx}x$\\
$w_{min,cash} = 1 - \dfrac{\mu_P-r_f}{\bar\mu^tw_T}$\\
$w_{min} = (1-w_{min,cash})w_T$\\
$\sigma_{min} = \sqrt{w_{min}^t\Sigma_Rw_{min}}$
\end{minipage}}
\end{table}

Note that the asset allocation formulas (9.13--9.14) and (9.15--9.16) will be derived in
section 9.3.
